\section{Reflection on Semester One}

Semester One of the TAC Industries project focused on establishing a strong architectural and conceptual foundation rather than delivering a feature-complete system. This approach proved effective for a security-critical project, as it allowed key design decisions to be explored, justified, and validated before committing to extensive implementation work.

A primary outcome of this phase was a clear understanding of how identity, access control, and network segmentation must interact to form a cohesive Zero Trust architecture. Designing the system from an architecture-first perspective highlighted the importance of separating reachability from authorisation and enforcing separation of duties across both infrastructure and application layers.

Limited proof-of-concept activities reinforced the feasibility of the proposed design while also revealing practical constraints and trade-offs. These insights informed refinements to trust boundaries, access flows, and monitoring assumptions. Overall, Semester One successfully achieved its goal of producing a realistic, secure, and ethically grounded design that is suitable for further development.

\section{Planned Enhancements for Semester Two}

Semester Two will focus on extending the architectural foundation established in Semester One into a more complete and operational system. Planned enhancements include deeper implementation of application-level access control, enforcement of fine-grained role, attribute-based policies, an integrated web app for managing the cases created and expansion of auditing and monitoring capabilities.

Additional work will be undertaken to harden the infrastructure, improve automation, and refine identity security controls. This includes enhanced privileged access management, stronger authentication enforcement, and improved validation of security events. Greater emphasis will also be placed on documenting implementation outcomes and evaluating the effectiveness of security controls through testing and observation.

By building upon the validated design produced in Semester One, Semester Two aims to transition the TAC Industries platform from a conceptual and partially validated architecture into a more fully realised and demonstrable system.